\documentclass[10pt,a4paper]{amsart}
\usepackage[margin=19mm]{geometry}
\usepackage{amsmath,amssymb,mathtools}
\usepackage[scr=rsfs,cal=euler]{mathalfa}
\usepackage{xfrac}
\usepackage[unicode,hidelinks,bookmarks=false]{hyperref}
\newcommand{\ie}{i.e.\ }
\newcommand{\cf}{cf.\ }
\newcommand{\resp}{respectively\ }
\newcommand{\Subsection}{\subsection}
\providecommand{\nobreakdash}{\nobreak\hbox{-}}
\setlength{\emergencystretch}{2em}
\multlinegap0pt
\newtheorem{theorem1}{Theorem1}
\usepackage[shortlabels]{enumitem}
\usepackage{amssymb}
\usepackage[all]{xy}
\usepackage{tikz}
\newtheorem{lemma}{Lemma}
\newtheorem{theorem}{Theorem}
\title{Optimal Quantization on Spherical Surfaces: \\ Continuous and Discrete Models -- A Beginner-Friendly Expository Study}
\author{Mrinal Kanti Roychowdhury}
\date{Source: arXiv:2511.05099v2 (CC BY 4.0)}
\begin{document}
\maketitle

\noindent\emph{Source and licence.} Optimal Quantization on Spherical Surfaces: \\ Continuous and Discrete Models -- A Beginner-Friendly Expository Study. Version arXiv:2511.05099v2 (CC BY 4.0), distributed by arXiv under the Creative Commons Attribution 4.0 International licence, http://creativecommons.org/licenses/by/4.0/, as recorded in the arXiv licence metadata for this exact version and shown on the abstract page https://arxiv.org/abs/2511.05099v2. Full licence notice: \texttt{licenses/arxiv-2511.05099-CC-BY-4.0.txt}.\par\smallskip

\noindent\textit{Editorial scope.} Counted unit: the article's theorem thm:discrete-arc (source0.tex lines 1164--1174). The article's complete original proof of it is delivered with it (source0.tex lines 1176--1178, one \texttt{proof} environment of 3 lines). No mathematics is written, repaired or re-derived: the statement and the proof are the article's own lines, in the article's own notation, and no retained line is substituted or re-flowed.  Retained uncounted as the source-local context the proof needs: lines 1129--1138; lines 1156--1158.  Nothing was omitted from any retained span, so the package carries no omission record.  No retained line contains a citation, so the wrapper carries no bibliography and no bibliography entry.  No retained line contains a figure environment, an image-inclusion command or drawing code, so no image is delivered and the package declares no asset dependency.  Editorial formatting only, in addition: an AMS \texttt{amsart} preamble that restates the article's own declarations and macros used by the retained lines, namely the environments \texttt{lemma}, \texttt{theorem} on separate counters, together with the standard packages those lines need.  The retained environments print in this document as Lemma 1, Theorem 1 in order of appearance, whereas the article prints the same items with the numbering of its own full document.  The complete native source of the article as distributed in the arXiv e-print for this exact version is bundled unchanged as \texttt{sources/188636/source0.tex}.
\begin{theorem1}[Uniform arc: structure and error]\label{thm:arc-continuous}
Let $P_A$ be uniform on a great circular arc $A$ of length $L$. For each $n\ge 1$:
\begin{enumerate}[label=(\roman*), leftmargin=1.5em]
  \item The optimal set of $n$-means is obtained by partitioning $A$ into $n$ adjacent sub-arcs of equal length $L/n$ and placing each representative at the midpoint of its sub-arc.
  \item The $n$th quantization error is
  \[
  V_n(P_A)\;=\;\frac{L^2}{12\,n^2}.
  \]
\end{enumerate}
\end{theorem1}

\begin{lemma}[One-block center is the (discrete) midpoint]\label{lem:block-mid}
Fix a contiguous block $B=\{i,\dots,j\}$ and consider $f(a)=\sum_{\ell=i}^j d_G^2(x_\ell,a)$ for $a$ constrained to $A$. Then any minimizer $a^\ast$ lies at the \emph{geodesic midpoint} of the endpoints of the block, and if the block has odd cardinality, $a^\ast$ is the central data point; if even, any point in the middle geodesic segment between the two central data points minimizes $f$.
\end{lemma}

\begin{theorem}[Discrete-uniform arc: optimality and error]\label{thm:discrete-arc}
Let $P_X$ be the discrete-uniform measure on $m$ equally spaced points on an arc $A$ of length $L$. For each $n\le m$:
\begin{enumerate}[label=(\roman*), leftmargin=1.5em]
  \item An optimal $n$-clustering is obtained by partitioning into $n$ contiguous blocks whose sizes differ by at most $1$; the associated block centers (as in Lemma~\ref{lem:block-mid}) form an optimal set of $n$-means.
  \item If $m$ is a multiple of $n$, so each block has $m/n$ points and span $L/n$, then
  \[
  V_n(P_X)\;=\;\frac{1}{m}\cdot n\sum_{k=-(m/n-1)/2}^{(m/n-1)/2}\!\!\!\!\!\! \bigg(\frac{k\,L}{m}\bigg)^2,
  \]
  which, for large $m$, converges to the continuous value $L^2/(12n^2)$.
\end{enumerate}
\end{theorem}

\begin{proof}
(i) By convexity of the one-block objective and the exchange argument from Theorem~\ref{thm:arc-continuous}, blocks must be contiguous and as balanced as possible. (ii) With uniform spacing $\Delta=L/(m-1)$ (or approximately $L/m$ for large $m$), the sum of squared distances in a block centered at its midpoint is a symmetric discrete quadratic sum. Dividing by $m$ and summing over $n$ blocks yields the stated formula, which approaches the continuous integral $h^2/12$ per block with $h=L/n$ as $m\to\infty$.
\end{proof}

\end{document}
